% Section 4.1 -- Detection Evaluation Protocol
% Standalone, compilable draft (FPT-template article class + booktabs).
% Folded into the thesis main document later via \input.
%
% Honesty-phrasing rules applied (personal_agenda P1):
%   - "mean IoU" always qualified as "of matched detections"
%   - recall denominator stated explicitly (>=10 points surviving preprocessing)
%   - evaluation is leave-one-sequence-out (LOSO): each sequence is scored by a
%     classifier that never saw it in training or selection; seq 08 is the reported
%     operating point (leakage-free fold-08 confirms the final seq-08 numbers)
%   - measured vs inferred evidence tiers marked in prose
%   - every number traceable to a finding / output artifact
%
\chapter{Detection Evaluation}
\label{ch:detection-eval}

\section{Evaluation protocol}
\label{sec:eval-protocol}

Every quantitative claim about detection in this thesis is produced by a single
evaluation script (\texttt{src/evaluate.py}) under one fixed protocol. Because
the results, ablation, and recall-analysis chapters all report numbers from this
protocol, we state it in full here and refer back to it rather than restating it
per table. The protocol was frozen before the write-up began; no configuration,
checkpoint, or metric definition described below was altered to produce a
reported number.

\subsection{Sequences, roles, and leave-one-sequence-out cross-validation}
\label{sec:eval-sequences}

We evaluate on the eleven labelled SemanticKITTI odometry sequences
(\numrange{00}{10}), whose per-point semantic and instance labels serve as
ground truth. The official test sequences \numrange{11}{21} carry \emph{no}
public labels, so no ground-truth detection metric can be computed on them; this
is a property of the dataset, not a choice, and it is why no labelled sequence
can be reserved as a conventional held-out test set.

\paragraph{Classifier evaluation is leakage-free by construction.}
Rather than report detection on the sequence the classifier was tuned against, we
evaluate the binary classifier under \emph{leave-one-sequence-out} (LOSO)
cross-validation over all eleven sequences. For each held-out sequence~$K$ a
classifier is trained on the other ten sequences and then run on the full
sequence~$K$; the held-out sequence contributes \emph{zero} clusters to that
fold's training data and is never used to select the checkpoint. Checkpoint
selection is removed entirely: every fold trains for a fixed \num{15}-epoch
budget and the last-epoch checkpoint is used (historically the best epoch was
\num{13}/\num{15}, so last~$\approx$~best), so the held-out sequence is untouched
during both training and model selection. Every per-sequence number in the
per-fold table below is therefore leakage-free with respect to the classifier:
the classifier scoring sequence~$K$ has never seen sequence~$K$. Because the
Stage~B clusters are stored with a sequence prefix in their filenames, the folds
are a re-partition of the existing mined data and require no re-mining.

\paragraph{Two sequences carry additional, distinct roles.}
\begin{itemize}
  \item \textbf{Sequence 08 (reported operating point; \num{4071} frames).} The
    largest labelled sequence is the operating point at which the pipeline is
    reported in detail: the stage ablations (Section~\ref{sec:results-ablations}),
    the recall root-cause analysis (Chapter~\ref{ch:recall-bottleneck}), the
    IoU-threshold sensitivity below, and the completion evaluation
    (Chapter~\ref{ch:completion-eval}) are all computed on the sequence-08 output
    of the \emph{final} classifier. Sequence~08 was never used to tune the
    classical stages. The final classifier was historically selected on
    sequence~08, so its sequence-08 numbers could in principle be optimistic; the
    LOSO fold-08 classifier (trained on \numrange{00}{07},\,\numrange{09}{10},
    last-epoch, never selected on 08) settles this empirically below.
  \item \textbf{Sequence 00 (completion held-out; 100 frames).} Sequence~00 is
    the held-out sequence for the \emph{completion} replication
    (Chapter~\ref{ch:completion-eval}), where its being unseen by the completion
    model is what matters. Its classical pipeline parameters
    (\texttt{PIPELINE\_CONFIG}) were historically hand-tuned on it, so the
    seq-00 detection fold retains classical-stage familiarity even though its
    classifier is still properly held out; we isolate that residual in the
    aggregate below (dropping fold-00 moves pooled recall by \num{0.022}).
\end{itemize}

\subsection{Ground-truth targets and the eligibility rule}
\label{sec:eval-targets}

The pipeline is a car detector, so the evaluation counts only the two
car-bearing SemanticKITTI classes as valid ground-truth objects:
\emph{car} (class~10) and \emph{moving-car} (class~252); we refer to this target
set as \emph{supported-vehicles}. A ground-truth car \emph{instance} is a
maximal set of points sharing one instance id within these classes.

An instance is counted in the evaluation only if at least \textbf{10 of its
points survive preprocessing} --- the $z$-filter, voxel downsampling, statistical
denoising, and ground removal that every scan passes through before clustering.
This is the recall denominator: \emph{recall is computed against ground-truth car
instances with $\geq 10$ points surviving preprocessing}, not against every
annotated car. The rule exists because an instance reduced to a handful of points
by occlusion or range cannot be segmented by any clustering method, so including
it would measure sensor sparsity rather than the pipeline.

This choice narrows the denominator, and honesty requires quantifying by how
much. On a stride-20 sample of sequence~08 (204 of \num{4071} frames, drawn
uniformly across the drive), of all annotated car instances
(\num{2332} pooled), \textbf{25.5\% are excluded} by the surviving-points rule;
measured instead against instances that have $\geq 10$ points in the \emph{raw}
scan (\num{2033} pooled), \textbf{14.6\% are excluded} by preprocessing.\footnote{%
These are two distinct $\geq 10$-point rules: one thresholds \emph{raw} label
points (a distinct-car count), the other thresholds points \emph{surviving
preprocessing} (the eval denominator). They are not interchangeable.}
Because an excluded instance can never be matched, the sequence-08 operating-point
recall of \num{0.730} against the eligible denominator corresponds, by inference, to a
recall of \emph{approximately} \num{0.54} against \emph{all} annotated cars
(the eligible-to-annotated survival ratio on this sample is
$\num{1737}/\num{2332}=\num{0.745}$; the ratio is a stride-20 estimate, so this
figure is inferred, not separately measured). We report the eligible-denominator
recall as the headline and disclose the all-annotated figure alongside it.

\subsection{Matching and metrics}
\label{sec:eval-matching}

Detections and ground-truth instances are compared at the \emph{point level}.
For a detection mask $D$ and an instance mask $G$ over the surviving points, the
overlap is the intersection-over-union
\[
  \mathrm{IoU}(D,G) \;=\; \frac{|D \cap G|}{|D \cup G|}.
\]
All detection--instance pairs with $\mathrm{IoU} \geq \tau$ are collected, sorted
by descending IoU, and assigned \emph{greedily and one-to-one}: the highest-IoU
pair is matched first, and once a detection or an instance is used it cannot be
matched again. This yields, per frame:
\begin{itemize}
  \item a \textbf{true positive} (TP) for each matched pair,
  \item a \textbf{false positive} (FP) for each unmatched detection,
  \item a \textbf{false negative} (FN) for each unmatched eligible instance.
\end{itemize}
Precision, recall, and $F_1$ are then computed from TP, FP, and FN
\emph{micro-averaged} over the sequence: counts are pooled across all frames and
the ratios taken once, so denser frames contribute proportionally more. We also
report the \textbf{mean IoU of matched detections} --- the mean of the IoU values
of the accepted TP pairs only. This is a quality measure of the matches that were
made; it says nothing about missed cars and must not be read as an overall
segmentation score.

Track-level filtering is enabled by default for the headline numbers: an offline
centroid tracker links detections across frames, and only tracks meeting a
minimum length and class-consistency requirement are accepted. This is a
post-hoc, offline step (it uses the whole sequence), consistent with the
offline scope of the system; its effect is isolated as an ablation row in the
ablation section (Section~\ref{sec:results-ablations}).

\subsection{Operating point}
\label{sec:eval-operating-point}

Table~\ref{tab:headline} gives the reported operating point on the full
sequence~08 under this protocol, produced by the final classifier (the one the
completion pipeline actually runs with). This is the operating point the rest of
this chapter and the analysis chapters decompose; the leakage-free
cross-validation that confirms it follows in Section~\ref{sec:eval-loso}.

\begin{table}[H]
  \centering
  \caption{Reported detection operating point under the frozen protocol, using
    the final classifier. Sequence~08 is the largest labelled sequence and was
    never used to tune the classical stages; sequence~00 is a completion held-out
    reference whose classifier saw seq-00 clusters in training (so its detection
    recall reads optimistically). Mean IoU is the mean over matched detections
    only. Section~\ref{sec:eval-loso} re-scores every sequence with a classifier
    that never saw it.}
  \label{tab:headline}
  \setlength{\tabcolsep}{5pt}
  \begin{tabular}{lccccc}
    \toprule
    Sequence & Precision & Recall & $F_1$ & Mean IoU & TP / FP / FN \\
    \midrule
    08 (full, \num{4071} fr.) & 0.905 & 0.730 & 0.808 & 0.912 & \num{25478} / \num{2676} / \num{9444} \\
    00 (100 fr., reference)   & 0.967 & 0.777 & 0.862 & 0.962 & \num{1296} / \num{44} / \num{371} \\
    \bottomrule
  \end{tabular}
\end{table}

The pipeline is precision-saturated and recall-limited: false positives are flat
at roughly one per frame, while recall is set by how many eligible cars the
clustering stage recovers as single objects. The remaining sections of this
chapter report the ablations that decompose this operating point, and Chapter~\ref{ch:recall-bottleneck}
root-causes the recall shortfall.

\subsection{Leave-one-sequence-out cross-validation}
\label{sec:eval-loso}

The operating point above is reported on sequence~08, which is also the sequence
the final classifier was historically selected on. To show that this selection
does not inflate the result, and to characterise the per-sequence recall spread
that a single reporting sequence cannot, we ran the LOSO protocol of
Section~\ref{sec:eval-sequences}: eleven folds, each training a classifier on ten
sequences and evaluating the full held-out eleventh, so every fold's classifier
is blind to the sequence it scores. Table~\ref{tab:loso} reports all eleven folds,
the pooled micro-average (counts summed across folds, comparable to the single
operating-point number), and the macro mean~$\pm$~standard deviation across folds.

\begin{table}[H]
  \centering
  \caption{Leave-one-sequence-out cross-validation over all eleven labelled
    sequences (full sequences, $\tau=0.3$, final configuration). Each fold's
    classifier trained on the other ten sequences and was never selected on the
    held-out one, so every row is leakage-free with respect to the classifier.
    ``Pooled'' sums TP/FP/FN across folds; the excl.-seq00 pool isolates the
    residual seq-00 classical-tuning familiarity (Section~\ref{sec:eval-sequences}).}
  \label{tab:loso}
  \setlength{\tabcolsep}{3pt}
  \footnotesize
  \begin{tabular}{lccccc}
    \toprule
    Held-out seq & Precision & Recall & $F_1$ & Mean IoU & TP / FP / FN \\
    \midrule
    00 (\num{4541} fr.) & 0.913 & 0.761 & 0.830 & 0.932 & \num{42763} / \num{4079} / \num{13427} \\
    01 (\num{1101} fr.) & 0.951 & 0.336 & 0.496 & 0.953 & \num{539} / \num{28} / \num{1067} \\
    02 (\num{4661} fr.) & 0.793 & 0.700 & 0.743 & 0.903 & \num{10919} / \num{2857} / \num{4685} \\
    03 (\num{801} fr.)  & 0.733 & 0.623 & 0.673 & 0.904 & \num{1527} / \num{556} / \num{926} \\
    04 (\num{271} fr.)  & 0.797 & 0.501 & 0.615 & 0.835 & \num{529} / \num{135} / \num{527} \\
    05 (\num{2761} fr.) & 0.861 & 0.731 & 0.791 & 0.910 & \num{10671} / \num{1723} / \num{3927} \\
    06 (\num{1101} fr.) & 0.861 & 0.637 & 0.732 & 0.906 & \num{6374} / \num{1031} / \num{3637} \\
    07 (\num{1101} fr.) & 0.905 & 0.729 & 0.808 & 0.929 & \num{11131} / \num{1162} / \num{4137} \\
    08 (\num{4071} fr.) & 0.881 & 0.737 & 0.803 & 0.912 & \num{25733} / \num{3477} / \num{9189} \\
    09 (\num{1591} fr.) & 0.824 & 0.646 & 0.724 & 0.911 & \num{5141} / \num{1097} / \num{2813} \\
    10 (\num{1201} fr.) & 0.725 & 0.559 & 0.632 & 0.871 & \num{1832} / \num{694} / \num{1444} \\
    \midrule
    \textbf{Pooled (all 11)} & \textbf{0.874} & \textbf{0.719} & \textbf{0.789} & \textbf{0.919} & \num{117159} / \num{16839} / \num{45779} \\
    Pooled (excl.\ seq~00)   & 0.854 & 0.697 & 0.767 & 0.911 & \num{74396} / \num{12760} / \num{32352} \\
    Macro mean $\pm$ std     & $0.840{\pm}0.070$ & $0.633{\pm}0.121$ & $0.713{\pm}0.097$ & $0.906{\pm}0.030$ & --- \\
    \bottomrule
  \end{tabular}
\end{table}

Three things follow. First, \textbf{model selection did not inflate the reported
recall}. The leakage-free fold-08 classifier reaches recall \num{0.737} against
the sequence-08 operating point's \num{0.730} --- if anything marginally higher,
not lower --- so the recall headline is not a product of having selected the
classifier on the reporting sequence. The model-selection optimism that does
exist is confined to precision and is small: fold-08 precision is \num{0.881}
versus the final \num{0.905} (an extra \num{801} false positives across
\num{4071} frames), exactly the two-point, precision-only exposure the stage
decomposition predicts (Section~\ref{sec:disc-eval-validity}). The two
sequence-08 rows are not two separate experiments but \emph{one classifier under
two checkpoint-selection rules}: the final model and the fold-08 model share the
same architecture, the same ten training sequences (\numrange{00}{07},
\numrange{09}{10}), and the same from-scratch \num{15}-epoch regime, and differ
only in which epoch is frozen --- the final model takes the epoch that scored
best on the sequence-08 validation set, the fold-08 model takes the last epoch and
never inspects sequence~08 at all. The \num{0.730}~$\rightarrow$~\num{0.737}
(recall) and \num{0.905}~$\rightarrow$~\num{0.881} (precision) deltas therefore
isolate the effect of that selection: selecting on sequence~08 bought about two
precision points and no recall. The isolation is not perfect --- because
sequence~08 cannot serve as the fold's loss-monitoring split, the fold-08 model
trains on \SI{95}{\percent} of the final model's clusters (a \num{5}-percent
monitoring slice is set aside) --- but that confound removes training data from
the leakage-free model, which would depress its recall, not inflate it, so it only
strengthens the ``recall was not inflated'' reading. Second, the
\textbf{pooled cross-validated operating point} over all eleven sequences ---
precision \num{0.874}, recall \num{0.719}, $F_1$~\num{0.789}, mean IoU
\num{0.919} --- sits within roughly one point of the sequence-08 operating point
on every metric, so the reported numbers are representative of the labelled set
and not specific to sequence~08. Third, the \textbf{per-sequence spread} is now
explicit: recall ranges from \num{0.336} on sequence~01 (a sparse highway drive)
to \num{0.761} on sequence~00 (dense urban), with a macro mean of
\num{0.633}~$\pm$~\num{0.121}; precision is far steadier
(\num{0.840}~$\pm$~\num{0.070}). The gap between the pooled recall (\num{0.719})
and the macro mean (\num{0.633}) is driven by the smallest, sparsest sequences
(seq~01 and seq~04, the latter only \num{271} frames), which the count-weighted
pool down-weights and the unweighted mean does not; both figures are reported so
neither hides the spread. Finally, the residual seq-00 classical-tuning exposure
is bounded: dropping the seq-00 fold lowers pooled recall only from \num{0.719}
to \num{0.697} (a \num{0.022} delta), so the historical seq-00 pipeline tuning
does not drive the aggregate.

\subsection{Choice of IoU threshold and its sensitivity}
\label{sec:eval-iou-threshold}

All headline numbers use $\tau = 0.3$. A point-level threshold of 0.3 is
permissive relative to the bounding-box 0.5 common in detection benchmarks, so
it is fair to ask whether the results are an artifact of a lenient gate. They
are not, and the reason is visible in Table~\ref{tab:headline}: the mean IoU of
matched detections is \num{0.912} on sequence~08, far above either threshold, so
the matches are not marginal overlaps sitting just past $\tau$.

To make this concrete we re-ran the full sequence~08 evaluation at
$\tau \in \{0.25, 0.30, 0.50\}$ with everything else frozen
(Table~\ref{tab:iou-sensitivity}).

\begin{table}[H]
  \centering
  \caption{IoU-threshold sensitivity on the full sequence~08 (frozen
    configuration). The $\tau = 0.30$ row is the headline. $F_1$ moves by at most
    \num{2.6} points across the tested range (\num{0.811} at $\tau=0.25$ to
    \num{0.785} at $\tau=0.50$); mean IoU of matched detections rises at the
    stricter gate because it keeps only the better-overlapping matches.}
  \label{tab:iou-sensitivity}
  \begin{tabular}{cccccc}
    \toprule
    $\tau$ & Precision & Recall & $F_1$ & Mean IoU & TP / FP / FN \\
    \midrule
    0.25 & 0.908 & 0.732 & 0.811 & 0.910 & \num{25576} / \num{2578} / \num{9346} \\
    0.30 & 0.905 & 0.730 & 0.808 & 0.912 & \num{25478} / \num{2676} / \num{9444} \\
    0.50 & 0.879 & 0.709 & 0.785 & 0.927 & \num{24744} / \num{3410} / \num{10178} \\
    \bottomrule
  \end{tabular}
\end{table}

$F_1$ falls only \num{2.3} points from the headline $\tau = 0.30$ to the strict
$\tau = 0.50$ (\num{0.808} to \num{0.785}), and \num{2.6} points across the full
tested range from the permissive $\tau = 0.25$ (\num{0.811}), so the reported
quality is stable under the choice of threshold and does not depend on the lenient
value. We had
pre-registered that a threshold-sensitive result would be \emph{reported as a
finding rather than tuned away}; it was not sensitive.
